% SPDX-FileCopyrightText: 2016-2026 Lari Natri <lari.natri@iki.fi>
% SPDX-License-Identifier: GPL-3.0-or-later
%
% Chapter corner marks.
% This file is part of the 'ulsbs' package.
%
% Requires TikZ and style defaults.
% Loaded with @ a letter, before header-style setup; no song-lifecycle dependency.
% Owns mark dimensions, dashed flag and cached PGF geometry. The surrounding
% style updates chapter color/dashed state; this file installs no chapter hooks.
% Header/footer callers draw overlays for the page being built, not while song
% source is captured. Default geometry is computed at load time.

\ProvidesFile{ulsbs-chapter-marks.code.tex}[ULSBS chapter corner marks]

% Setup the colored chapter corner mark for all pages within ulMainChapters.
%
% \ulChapterCornerMarkOffset should be at least 2mm to leave room for imprecise
% cutting of the page. \ulChapterCornerMarkOffset + \ulChapterCornerMarkWidth should
% be less than 4 mm, so that the whole mark is left outside of printer margins
% when printing on home printer to a paper that requires no cutting.
\newdimen\ulChapterCornerMarkOffset\ulChapterCornerMarkOffset=2.5mm%
\newdimen\ulChapterCornerMarkWidth\ulChapterCornerMarkWidth=.954915mm% 2.5 * 0.618034^2 :)
\newdimen\ulChapterCornerMarkLength\ulChapterCornerMarkLength=\dimexpr\ulTopMargin+\ulHeadHeight-\ulChapterCornerMarkOffset\relax%
% Gap count: 3 gaps -> 7 segments -> gaps are segments 2, 4 and 6.
\newcount\ulChapterCornerMarkDashCount\ulChapterCornerMarkDashCount=3%

% Set this to true to have dashed chapter corner marks. The surrounding style's
% chapter setup resets it to false; this module only reads the flag when drawing.
\ulsbs@newbool\ifulChapterCornerMarkDashed\ulsbs@chapcornermarkdashedtrue\ulsbs@chapcornermarkdashedfalse\ulChapterCornerMarkDashed\ulsbs@chapcornermarkdashedfalse

\newdimen\ulsbs@ccm@halfw
\newdimen\ulsbs@ccm@arm
\newdimen\ulsbs@ccm@reach
\newdimen\ulsbs@ccm@start
\newdimen\ulsbs@ccm@end
\newcount\ulsbs@ccm@segments
\newcount\ulsbs@ccm@segment

% Boundary #1 on either leg. Compute from the total so rounding never leaves
% a stray bit of dash or shortens the last segment.
\newcommand*\ulsbs@ccm@boundary[1]{\dimexpr\ulsbs@ccm@halfw
  +\ulsbs@ccm@arm*\numexpr#1\relax/\ulsbs@ccm@segments\relax}

% Build a single filled path with its elbow at the origin and legs along +x/+y.
% In dashed mode, the elbow includes the first dash of both legs; the remaining
% dashes are disconnected subpaths. No edges meet between separate paint calls.
\newcommand*\ulsbs@ccm@build{%
  \pgftransformreset
  \ulsbs@ccm@halfw=\dimexpr\ulChapterCornerMarkWidth/2\relax
  \ulsbs@ccm@reach=\ulChapterCornerMarkLength
  % A mark no longer than half its width consists of just the elbow.
  \ifdim\ulsbs@ccm@reach<\ulsbs@ccm@halfw
    \ulsbs@ccm@reach=\ulsbs@ccm@halfw
  \fi
  \ulsbs@ccm@segments=1
  \ifulChapterCornerMarkDashed
    \ifnum\ulChapterCornerMarkDashCount>0
      \ulsbs@ccm@segments=\numexpr2*\ulChapterCornerMarkDashCount+1\relax
    \fi
  \fi
  \ulsbs@ccm@arm=\dimexpr\ulsbs@ccm@reach-\ulsbs@ccm@halfw\relax
  \ulsbs@ccm@reach=\ulsbs@ccm@boundary{1}%
  \edef\ulsbs@ccm@h{\the\ulsbs@ccm@halfw}%
  \edef\ulsbs@ccm@r{\the\ulsbs@ccm@reach}%
  \pgfpathqmoveto{-\ulsbs@ccm@h}{-\ulsbs@ccm@h}%
  \pgfpathqlineto{ \ulsbs@ccm@r}{-\ulsbs@ccm@h}%
  \pgfpathqlineto{ \ulsbs@ccm@r}{ \ulsbs@ccm@h}%
  \pgfpathqlineto{ \ulsbs@ccm@h}{ \ulsbs@ccm@h}%
  \pgfpathqlineto{ \ulsbs@ccm@h}{ \ulsbs@ccm@r}%
  \pgfpathqlineto{-\ulsbs@ccm@h}{ \ulsbs@ccm@r}%
  \pgfpathclose
  \ifdim\ulsbs@ccm@arm>0pt
    \ulsbs@ccm@segment=2
    \@whilenum\ulsbs@ccm@segment<\ulsbs@ccm@segments\do{%
      \ulsbs@ccm@start=\ulsbs@ccm@boundary{\ulsbs@ccm@segment}%
      \ulsbs@ccm@end=\ulsbs@ccm@boundary{\ulsbs@ccm@segment+1}%
      \pgfpathrectanglecorners{\pgfqpoint{\ulsbs@ccm@start}{-\ulsbs@ccm@halfw}}%
        {\pgfqpoint{\ulsbs@ccm@end}{\ulsbs@ccm@halfw}}%
      \pgfpathrectanglecorners{\pgfqpoint{-\ulsbs@ccm@halfw}{\ulsbs@ccm@start}}%
        {\pgfqpoint{\ulsbs@ccm@halfw}{\ulsbs@ccm@end}}%
      \advance\ulsbs@ccm@segment by 2
    }%
  \fi
  \pgfusepathqfill
}

% Keep the most recent geometry of each style. Offset, orientation and chapter
% color are applied at use time, so neither page parity nor chapter changes
% require rebuilding the object.
\newcommand*\ulsbs@ccm@use[1]{%
  \edef\ulsbs@ccm@key{\the\ulChapterCornerMarkWidth/\the\ulChapterCornerMarkLength
    \ifulChapterCornerMarkDashed/\the\ulChapterCornerMarkDashCount\fi}%
  \expandafter\ifx\csname ulsbs@ccm@key@#1\endcsname\ulsbs@ccm@key
  \else
    \pgfdefobject{ulsbs@ccm@object@#1}%
      {\pgfqpoint{-.5\ulChapterCornerMarkWidth}{-.5\ulChapterCornerMarkWidth}}%
      {\pgfqpoint{\ulChapterCornerMarkLength}{\ulChapterCornerMarkLength}}%
      {\ulsbs@ccm@build}%
    \expandafter\xdef\csname ulsbs@ccm@key@#1\endcsname{\ulsbs@ccm@key}%
  \fi
  \pgfuseobject{ulsbs@ccm@object@#1}%
}

% Draws one chapter corner mark inside the caller's TikZ picture.
% #1 = anchor (north east / north west / south east / south west)
% #2 = x sign: + or -
% #3 = y sign: + or -
\newcommand*\ulChapterCornerMark[3]{%
  \pgfscope
    \pgftransformshift{\pgfpointadd{\pgfpointanchor{current page}{#1}}%
      {\pgfqpoint{#2\ulChapterCornerMarkOffset}{#3\ulChapterCornerMarkOffset}}}%
    \edef\ulsbs@ccm@xscale{\number\numexpr#21\relax}%
    \edef\ulsbs@ccm@yscale{\number\numexpr#31\relax}%
    \pgftransformcm{\ulsbs@ccm@xscale}{0}{0}{\ulsbs@ccm@yscale}{\pgfpointorigin}%
    \pgflowlevelsynccm
    \pgfsetfillcolor{\ulsbs@currentchaptercolor}%
    \ifulChapterCornerMarkDashed
      \ulsbs@ccm@use{dashed}%
    \else
      \ulsbs@ccm@use{solid}%
    \fi
  \endpgfscope
}

% Creates chapter corner marks on the outer corners of the page
% #1 = which side is outer?  R (odd pages) or L (even pages)
% Emits a zero-layout-space remembered TikZ overlay only when chapter colors are
% enabled. Caller supplies page parity; this macro does not advance page state.
\newcommand*\ulChapterCornerMarksOuter[1]{%
  \ifulUseChapterColors
    \begin{tikzpicture}[overlay,remember picture]
      \ifx#1R
        \ulChapterCornerMark{north east}{-}{-}%
        \ulChapterCornerMark{south east}{-}{+}%
      \else
        \ulChapterCornerMark{north west}{+}{-}%
        \ulChapterCornerMark{south west}{+}{+}%
      \fi
    \end{tikzpicture}%
  \fi
}
